% Figure: the Rankine-Hugoniot curve versus the isentrope through the
% same initial state, in the pressure-specific-volume plane. The
% Hugoniot lies above the isentrope on the compression branch (a shock
% reaches a higher pressure for the same compression because it raises
% entropy). The Rayleigh line connects initial and final states.
% Use v = specific volume (normalised v1=1, p1=1). gamma=1.4.
% Hugoniot: p2/p1 = ( (g+1) - (g-1)(v2/v1) ) / ( (g+1)(v2/v1) - (g-1) )
%   g=1.4 -> num = 2.4 - 0.4*v ; den = 2.4*v - 0.4   (v = v2/v1 in [0.2,1])
% Isentrope: p2/p1 = (v2/v1)^(-g) = v^(-1.4)
\begin{figure}[htbp]
\centering
\begin{tikzpicture}
\begin{axis}[
  width=11cm, height=8cm,
  xlabel={specific volume ratio $v_2/v_1$},
  ylabel={pressure ratio $p_2/p_1$},
  xmin=0.15, xmax=1.6, ymin=0, ymax=9,
  axis lines=left,
  grid=major, grid style={enggray!20},
  tick label style={font=\scriptsize},
  label style={font=\small},
  legend style={font=\scriptsize, at={(0.97,0.97)}, anchor=north east},
  legend cell align=left,
]
% Hugoniot curve (compression branch v<1 and a little of v>1)
\addplot[engred, very thick, domain=0.18:1.55, samples=160]
  { (2.4 - 0.4*x) / (2.4*x - 0.4) };
\addlegendentry{Hugoniot (shock)}
% Isentrope through (1,1)
\addplot[engblue, very thick, domain=0.18:1.55, samples=160]
  { x^(-1.4) };
\addlegendentry{isentrope}
% Rayleigh line from initial state (1,1) to a final state at v=0.4
% p at v=0.4 on Hugoniot: (2.4-0.16)/(0.96-0.4)=2.24/0.56=4.0
\addplot[englime, thick, dashed] coordinates {(1,1) (0.4,4.0)};
\addlegendentry{Rayleigh line}
% initial state marker
\fill[engblack] (axis cs:1,1) circle (1.6pt);
\node[engblack, font=\scriptsize, anchor=north west] at (axis cs:1.02,1.0)
  {$(1,1)$ upstream};
% final state marker
\fill[engblack] (axis cs:0.4,4.0) circle (1.6pt);
\node[engblack, font=\scriptsize, anchor=south west] at (axis cs:0.42,4.0)
  {downstream};
% vertical asymptote of Hugoniot at v=(g-1)/(g+1)=0.1667
\draw[enggray, dotted] (axis cs:0.1667,0) -- (axis cs:0.1667,9);
\node[enggray, font=\scriptsize, anchor=south west, rotate=90]
  at (axis cs:0.1667,3.0) {$v_2/v_1\to\frac{\gamma-1}{\gamma+1}$};
\end{axis}
\end{tikzpicture}
\caption[Hugoniot curve versus the isentrope]{The Rankine-Hugoniot
curve (red) and the isentrope (blue) drawn through the same upstream
state $(v_2/v_1,p_2/p_1)=(1,1)$, in the pressure-specific-volume plane
for air ($\gamma=1.4$). Both pass through the upstream point and both
rise as the gas is compressed ($v_2/v_1<1$), but on the compression
branch the Hugoniot lies above the isentrope: a shock reaches a higher
pressure for the same volume compression because it raises the entropy,
turning ordered kinetic energy into heat. The gap between the two curves
is the irreversibility of the shock. The Hugoniot has a vertical
asymptote at $v_2/v_1=(\gamma-1)/(\gamma+1)=1/6$, the limiting density
ratio of a perfect-gas normal shock: no shock, however strong,
compresses the gas past six times its upstream density. The Rayleigh
line (green dashed) is the locus of states that conserve mass and
momentum across the shock; its intersection with the Hugoniot (which
adds energy conservation) fixes the downstream state, and its slope sets
the shock speed. The expansion branch ($v_2/v_1>1$) of the Hugoniot
would lower the entropy, so the second law forbids expansion shocks: a
supersonic flow expands isentropically through a Prandtl-Meyer fan, not
through a discontinuity.}
\label{fig:vol03ch13:hugoniot}
\end{figure}
